\documentclass[instructions.tex]{subfiles}
\begin{document}
 
\chapter{Du Deul et des châtimens dont Dieu Menace les vindicatifs.}

\primo Le duel est un des grands écueils de la noblesse et de la professions des armes. Le duel, disent-ils, est une loi de la guerre. Mais qui a établie cette loi ? Ce n'est pas Dieu, puisqu’il le défend comme une action exécrable. Ce n’est pas la religion, puisque l’Église refuse ses suffrages et même la sépulture à ceux qui perdent la vie dans ces détestables combats. Ce ne sont pas les souverains, puisqu’ils emploient leur autorité pour exterminer ce brutal exercice.
 
Disons plutôt que c’est une loi de l’enfer ; et que le démon n’a jamais rien inspiré de plus horrible. L’homme n’a point, de son fonds, des sentimens si barbares et si contraires à sa nature. Ils ne peuvent être inspirés que par Satan, que l’Écriture appelle homicide dès le commencement, et qui ne se plaît qu’à la dégradation et à la perte de l’homme. Non, le duel n’est point une loi de la guerre, mais un égarement, une fureur, une frénésie, un attentat digne de tous les supplices.
 
Mais si je refuse le duel, je serai déshonoré. Devez-vous craindre d’être déshonoré dans l’esprit de ceux qui ont étouffé tout sentiment d’humanité et de religion ? S’exposer à la mort, par devoir, dans un combat, pour son prince et sa patrie, est une action pleine de générosité ; mais braver la mort par gloire dans le duel, pour tomber en mourant entre les mains du Juge éternel, est une action si pleine de folie, qu’il n’y a rien qui montre plus l’aveuglement des hommes, que d’avoir attaché de l’honneur à une action si basse et si indigne. Ce n’est pas ici qu’on reconnoît la grandeur d’ame des nobles, mais plutôt leur foiblesse.
 
La vraie noblesse est dans les sentimens d’humanité et de générosité, et non pas dans les sentimens d’une extravagante vanité. Job étoit noble et puissant prince; il pensoit noblement lorsqu’il étoit disposé à tout sacrifier pour sauver son âme. Salomon, le plus éclairé des rois, faisoit paroître des sentimens de noblesse et de grandeur lorsqu’il disoit : Honorez votre âme selon son mérite, et conservez-la dans les sentimens de clémence et de paix. Quelle bassesse de cœur ! pour une parole échappée, pour un passe-droit, ôter la vie à un homme ! perdre une ame plus précieuse que les empires de la terre ! s’exposer soi-même à perdre son honneur, sa famille, sa vie et son ame !
 
Si je refuse le combat, je perdrai mes emplois. Eh bien ! perdez tout, s’il le faut, mais ne vous perdez pas vous-même. Sacrifiez un point d’honneur imaginaire, et ne sacrifiez pas l’ame de votre frère. Son sang criera contre vous bien plus fort que le sang d’Abel contre Caïn.
 
Si l’on vous attaque, si l’on en veut à votre vie, défendez-vous avec modération, il vous est permis ; mais il ne vous est jamais permis d’accepter un rendez-vous pour le duel, ni de le demander, ni d’y concourir. Faites paroître votre bravoure contre les ennemis de l’État, mais non pas contre vos confrères et vos amis.
 
Combien de fléaux du ciel n’attire pas sur les royaumes ce barbare exercice ! Oh ! que ceux qui, dans la profession des armes, ont en main l’autorité, qui ne s’en servent pas pour abolir ces énormes abus, et qui les excusent, seront sévérement punis par le Dieu des armées !
 
II. Écoutez les menaces que Dieu a prononcées contre le vindicatif : Qu’il soit au pouvoir des méchans, dit le Prophète, et que le démon soit à sa droite pour le perdre ; qu’il soit condamné au tribunal de son juge ; que ses prières soient changées en péchés ; que sa vie soit abrégée et qu’un autre usurpe ses emplois ; qu’il laisse une femme dans un triste veuvage ; que ses enfans soient vagabonds, réduits à la mendicité, et chassés de leurs habitations ; que les étrangers consument ses biens ; que son nom et sa mémoire périssent dans une seule génération; que les iniquités de ses pères soient toujours présentes aux yeux du Seigneur, et que les péchés de sa mère ne soient jamais effacés. O mon Dieu ! quelles menaces ! Sur qui doivent-elles tomber ? sur le vindicatif ; parce que, dit le Prophète, il a refusé de pardonner, de faire miséricorde\footnote{Pro eo quòd non est recordatus facere misericordiam. Ps. 108.} .
 
Pourquoi le Seigneur menace-t-il ici de faire revivre les iniquités des pères ? c’est parce que les parens transmettent souvent leurs sentimens de vengeance et de haine à leurs enfans. Le saint Prophète ajoute que le persécuteur, le vindicatif, sera privé des bénédictions du ciel; qu’il s’est revêtu de la malédiction de Dieu; que cette malédiction pénétrera jusqu’à la moëlle de ses os, et ne le quittera point.
 
O combien de maisons désolées, de familles éteintes, d’ames perdues par la vengeance ! combien d’anathêmes et de malheurs n’attirent pas sur les hommes la haine et le peu de charité qu’ils ont les uns pour les autres !
 
\section{Résolutions}
 
\primo Je ne me vengerai désormais plus de mon ennemi, ni par mes discours, ni par mes actions. 

\secundo Je ferai même tous mes efforts pour étouffer dans mon cœur les pensées haineuses, les désirs du mal, la complaisance dans celui que j’entendrai dire, ou que j’apprendrai être arrivé.

\tertio A plus forte raison aurai-je toute ma vie une vive horreur du duel. 

\prayer{Mon Dieu, faites-moi la grâce d’arracher de mon cœur, dès que je m’en apercevrai, tout ce qui sentirait la haine et la vengeance.}
 
\end{document}
